% main.tex -- minimal wrapper for the Method and Experiments draft.
% Content source: ../METHOD_NEURIPS.md (English body + 中文要點 per subsection), converted 2026-10-17.
% Corrections of 2026-10-17 (../PROBLEMS_AND_SOLUTIONS.md, 「更正 2026-10-17」) are applied in the .tex wording;
% each place is marked with a "% NOTE:" comment, and the Markdown draft keeps the old wording with a 更正 mark.
% Project rules: one sentence per line; English only in typeset text; 中文要點 as % comments under each heading;
% every open item uses \tbd{...}; numbers are copied from the Markdown draft, not recomputed.
% Build (when a TeX system is available): pdflatex main && bibtex main && pdflatex main && pdflatex main
\documentclass{article}

% NeurIPS style: put the official neurips_2026.sty in this folder and uncomment the next line.
% It is not included here; do not substitute a hand-made copy.
% If the style loads natbib itself, keep the natbib line below option-free (or pass the style's nonatbib option).
% \usepackage[preprint]{neurips_2026}

\usepackage{amsmath}
\usepackage{amssymb}
\usepackage{booktabs}
\usepackage{natbib}
\usepackage{xcolor}
\usepackage{hyperref}

% Software without a year in our docs is cited by alias, so author-year citations do not print an empty year.
\defcitealias{opencc}{OpenCC}

% Open items stay visible in the PDF.
\newcommand{\tbd}[1]{\textcolor{red}{[TBD: #1]}}

% Chinese strings: #1 is the original (kept in the source, not typeset), #2 is what pdflatex prints (romanisation or gloss).
% With a CJK-capable build (e.g. xelatex + xeCJK) this can be redefined as \renewcommand{\zh}[2]{#1}.
\newcommand{\zh}[2]{#2}

\title{Fine-tuning a local language model to imitate one person's Traditional-Chinese chat style from a few hundred replies, with validated stylometric rulers as the training and acceptance signal}
\author{\tbd{authors}}
\date{}

% Main claim (METHOD_NEURIPS.md): an LLM can be fine-tuned, from a few hundred reply turns, to imitate a specific
% person's Traditional-Chinese (Taiwan colloquial) chat style, and this can be shown with rulers that were themselves
% validated on the same kind of text; the rulers are an intermediate step, not the contribution.
%
% Status (METHOD_NEURIPS.md, 2026-10-13): Method drafted from the validated parts of the pipeline (PIPELINE_SPEC.md);
% Experiments has two tiers (Tier 1 runnable now, Tier 2 planned). Numbers are from Experiments 1-4
% (PIPELINE_TECH_SURVEY.md, FORMAL_DISTANCE_EXPERIMENT.md); every number that depends on training is \tbd.
% Bibliography: refs.bib; entries carrying note = {... bibliography to verify} are incomplete in our docs.

\begin{document}

\maketitle

% Sections 1--2 (Introduction, Related work) are not drafted yet.
% Numbering starts at 3 so that section numbers match METHOD_NEURIPS.md.
\setcounter{section}{2}

% method.tex -- Section 3 (Method), converted from METHOD_NEURIPS.md §3 on 2026-10-17.
% One sentence per line. Each subsection's 中文要點 is the comment block directly under its heading.
% "% NOTE:" comments mark places where the wording differs from the Markdown draft, or where docs and code disagree.

\section{Method}
\label{sec:method}

\subsection{Problem setting}
\label{sec:setting}
% 中文要點:
% 單位是回合(同人、同日、間隔 ≤ 30 分);樣本 = 前 12 個回合 → T 的回合;
% 目標是生成的回覆被「分得出 T 與 S 的尺」分不出來,同時不離題。
% 資料三個特性:短、符號類內容多且隨時期變、只有兩位作者。

We are given the message log of one two-party chat between a \emph{target} person T and a \emph{partner} S, exported from LINE.
A \textbf{turn} is a maximal run of messages by one speaker on one day with gaps of at most 30 minutes between consecutive messages; a \textbf{segment} is a maximal run of turns with gaps of at most 6 hours.
The data form a sequence of segments, each a sequence of turns $(u_1, \dots, u_m)$ with speaker labels in $\{T, S\}$.

A \textbf{sample} is a pair $(c, y)$ where $y$ is a turn by T and $c$ is the context: the up-to-12 preceding turns of the same segment (both speakers), rendered as text.
The goal is a conditional generator $p_\theta(y \mid c)$ whose replies, under the conditions of the held-out period, cannot be told from T's own replies by rulers that \emph{can} tell T from S, while staying on topic.
% NOTE: the Markdown draft says the small training set "rules out the data regimes of prior persona work (thousands to
% tens of thousands of messages)". That conflicts with the 500-message IMPersona setting quoted in §4.1, so it is rephrased.
The training set is small: a few hundred samples (here about 700 after filtering), close to the smallest setting reported by IMPersona (500 messages; their largest uses 13K) \citep{shi2025impersona}.

Three properties of the data drive the design.
(i) Turns are short: the median turn has 16 characters, so a single turn carries little stylistic evidence.
(ii) The log is a tutoring chat, so a large and time-varying share of the characters (0.2--0.8 depending on the period) are code, formulas or symbols rather than conversation.
(iii) Only two authors exist, so classical multi-author stylometry has no reference population.

\subsection{Overview}
\label{sec:overview}
% 中文要點:
% 三階段:尺(先驗證、再校準)→ SFT LoRA → 尺排序的 DPO。驗收用另一組尺。

The method has three stages (\tbd{Figure 1: pipeline overview}).
\begin{enumerate}
\item \textbf{Rulers} (\S\ref{sec:rulers}): a small set of stylometric measures on bags of turns, each validated on T and S before use, with calibrated thresholds; plus interpretable feature tolerance bands.
\item \textbf{Supervised adaptation} (\S\ref{sec:sft}): LoRA fine-tuning \citep{hu2022lora} of a 4B base or instruct model on the samples, with loss only on T's turns.
\item \textbf{Ruler-guided preference optimisation} (\S\ref{sec:pref}): sample candidates for each training context, rank them with the calibrated ruler inside interpretable behaviour cells, and run DPO \citep{rafailov2023dpo} on the resulting pairs against a fixed reference, with a capped score and content constraints.
\end{enumerate}
Evaluation (\S\ref{sec:eval}) uses \emph{different} rulers from training: an independent discriminator trained on the test period, a model-free compression distance, feature-distribution checks and a blind judgement by T.

\subsection{Rulers: validated stylometric distances on bags of turns}
\label{sec:rulers}
% 中文要點:
% 主尺 = 字元 1–3-gram 作者 LM 的 LLR(袋內加總;門檻按袋大小校準、單回合分數用等張回歸校準);
% 交叉檢查 = NCD;Cosine Delta 暫列。
% 16 個可解釋特徵各有 T 的容忍帶(5–95 百分位),兩側都扣。
% (更正 2026-10-17:Markdown 原稿的中文要點寫 14 個;正確是 16 個,見 persona/rulers/features.py 的 FEATURES。)
% 三個有效性檢查:k=10 判定 ≥ 0.95、k=50 的 within<cross ≥ 95%、內容固定子集 AUC ≥ 0.90。
% 單回合 ≤ 0.80,所以逐候選的分數只能當排序訊號。

\paragraph{Preprocessing for the rulers.}
Turns are kept as text with newlines between messages and single placeholder characters for stickers, images and emoticons.
A turn is flagged \emph{symbol-like} if, after emoticon placeholders are inserted, it contains any of \verb!; { } ( ) = < > # [ ] _ \ / * + & | % ^ $ @! or a backtick, a URL, or a programming keyword.
Flagged turns are excluded from ruler fitting and scoring and are not used as targets; in contexts they are replaced by a placeholder.
This rule was chosen empirically (Exp.~4): without it, within-person variation in symbol content exceeds the between-person difference, and the five distances we tested all fail the validity check below; a density-based rule that keeps formulas and bracketed fragments is weaker.

\paragraph{Author language model (primary ruler).}
Let $q_T$ and $q_S$ be character 1--3-gram multinomial models with add-one smoothing fitted on T's and S's training turns \citep[cf.][]{peng2003language}.
For a text $x$ the score is the log-likelihood ratio $\ell(x) = \log q_T(x) - \log q_S(x)$ (plus the log prior ratio); for a bag $B$ of turns, $\ell(B) = \sum_{x \in B} \ell(x)$.
This is the likelihood-ratio statistic for ``written by T'' under the two fitted models (the analogue of a matched filter), and at each bag size tested in Exp.~3 it was more accurate than the formal distances we compared.
Because $\ell(B)$ is a sum, its decision threshold depends on $|B|$; using 0 is miscalibrated (0.877 accuracy at $k = 50$ against 1.000 calibrated).
We therefore tune a threshold per bag size on out-of-fold training bags from five contiguous date blocks, and map single-turn scores to calibrated log-odds by isotonic regression on the same out-of-fold scores, fitted with the pool-adjacent-violators (PAV) algorithm of \citet{ayer1955empirical}.
% NOTE (更正 2026-10-17, PROBLEMS_AND_SOLUTIONS.md S10): Brümmer & du Preez 2006 is a speaker-detection evaluation paper,
% not a forensic one, and C_llr is not a calibration loss by itself; the calibration part is C_llr - C_llr^min.
% The definition of C_llr^min should still be checked in the full text (PROBLEMS_AND_SOLUTIONS.md, last line).
We report the log-likelihood-ratio cost $C_{llr}$ introduced by \citet{brummer2006application} for the evaluation of speaker detection, which reflects discrimination and calibration together.
The calibration part alone is $C_{llr} - C_{llr}^{\min}$, where $C_{llr}^{\min}$ is the cost of the same scores after an optimal monotone (PAV) transformation; \tbd{report $C_{llr}^{\min}$ next to $C_{llr}$ for the single-turn calibration}.

\paragraph{Model-free cross-check.}
Normalised compression distance \citep[NCD;][]{cilibrasi2005clustering} is computed with zlib against ten reference bags of about 3,000 characters per author; the bag score is $\mathrm{NCD}(B, S) - \mathrm{NCD}(B, T)$.
NCD needs no vocabulary and was the most stable of our measures across periods (worst-block AUC 0.992 at $k = 50$), but its threshold does not transfer, so it is reported as AUC or pairwise accuracy only.

\paragraph{Provisional formal distance.}
Cosine Delta \citep{smith2011improving,evert2017understanding} is computed on the 300 most frequent 1--3-grams, with z-scores taken from 400 random training bags because there is no multi-author reference.
It passes the validity check after filtering but depends on the reference draw; it is kept for reporting and will be replaced by a multi-author reference when such data exist (Tier~2).

\paragraph{Interpretable features and tolerance bands.}
% NOTE: the feature list follows FEATURES in persona/rulers/features.py (16 entries). The Markdown draft's list omits
% sentence-final punctuation (end_punct) and calls the tilde feature an "ending"; the code counts messages that contain
% a full-width or half-width tilde anywhere.
Sixteen features are computed on the ruler population (non-symbol turns): sentence-final particles (two groups), full-width punctuation share, sentence-final punctuation, sticker rate, burst size, message length, messages containing a tilde (full-width or half-width), catchphrases (two groups), question rate, letter share, digit share, laughter, first person, and emoticons.
For each feature the 5th--95th percentile over 1,000 training bags of T ($k = 50$) is the \textbf{tolerance band}.
A bag of replies is penalised for leaving the band in either direction, so a generator that over-produces T's catchphrases is penalised as well as one that under-produces them.
Features whose bag-level T-vs-S AUC is below 0.6 would be reported but not weighted; on our data each of the sixteen is above 0.6.
The weighted feature distance uses weights proportional to each feature's discrimination (AUC minus 0.5) with $|z|$ capped at 5, so that near-constant rare features cannot dominate.

\paragraph{Validity checks (pre-registered, rerun whenever preprocessing changes).}
A measure may be used only if (a) on bags of $k = 10$ turns from the held-out period, calibrated attribution accuracy is at least 0.95; (b) on 1,000 random draws of training bags with $k = 50$, the cross-person distance exceeds both within-person distances in at least 95\% of draws, both when bags are drawn from the whole period and when all bags of a draw come from the same date block; (c) on the subset of short acknowledgement turns (content held fixed), out-of-fold single-turn AUC is at least 0.90.
% NOTE (更正 2026-10-17, PROBLEMS_AND_SOLUTIONS.md §4): Min-max similarity is the Ruzicka index, not Bray-Curtis.
% Bray-Curtis similarity is 2R/(1+R) for Min-max R, a monotone transform: rank-based results agree, values and thresholds do not.
% Cite Bray & Curtis (1957) only where Bray-Curtis itself is named.
% NOTE: the values in this paragraph are Exp. 4. The stage-1 rerun with the emoticon whitelist (reports/ruler_validation.md;
% handover.md §4 item 3) gives slightly different values for (b) at k = 50 and k = 25; pick one source before submission.
After filtering, the author LM reaches 1.000 [0.99, 1.00] at $k = 10$ and 0.944 on the content-controlled subset; Cosine Delta and NCD pass (b) at 0.96--0.97; Min-max similarity $\sum_i \min(x_i, y_i) / \sum_i \max(x_i, y_i)$ (the Ruzicka index) is at 0.94 (Exp.~4).
At $k = 25$ the distances fall short of (b) (0.73--0.89), so validity is claimed at $k = 50$ and attribution at $k = 10$--25, and the two are kept separate.

\paragraph{What the rulers cannot do.}
% NOTE: 0.80 is the tuned single-turn value of Exp. 4 (PIPELINE_SPEC.md: 0.77 / tuned 0.80 / AUC 0.885);
% handover.md §4 item 1 quotes tuned 0.81 (stage-1 rerun, 0.807 in reports/ruler_validation.md).
Single-turn attribution accuracy is 0.80 or below for each measure we tested, and 0.66--0.72 on seven-character turns.
Per-candidate scores are therefore used in training only as a noisy ranking signal (\S\ref{sec:pref}), not as point judgements.

\subsection{Supervised adaptation}
\label{sec:sft}
% 中文要點:
% 一個 prompt 函式,訓練推論共用;一則一行、每行有標記;目標含整個回合 + EOS(學連發);不用 chat template。
% LoRA r8 α16 全 linear;dev 選 epoch 後用全訓練期重訓;放金絲雀句稽核記憶。

\paragraph{Format.}
Each sample is rendered by one function used identically in training and inference: one message per line, prefixed with \zh{我:}{\textit{w\v{o}}:} (``me'', for T) or \zh{對方:}{\textit{du\`{\i}f\={a}ng}:} (``the other party'', for S); the target is all lines of T's turn followed by an end-of-sequence token.
The model therefore also learns \emph{whether to send another message} (burst size, median 2 for T vs 1 for S), which the rulers treat as a feature.
No chat template is used, for either the base or the instruct checkpoint, so the base-vs-instruct comparison is not confounded by templates and the prompt at inference cannot drift from the prompt in training.
Contexts are truncated from the oldest turn to 512 tokens; the loss is computed only on target tokens, and a unit test asserts that the first target token is included.

\paragraph{Fine-tuning.}
We use LoRA \citep{hu2022lora} with rank 8, $\alpha = 16$, dropout 0.05 on all linear projections of Qwen3.5-4B (base and instruct), bf16, learning rate $10^{-4}$ with cosine decay, effective batch 8, sequence length 768, and 3--5 epochs depending on the subset size, with the epoch chosen on a time-ordered development split (last 15\% of the training period) and the final model retrained on the full training period with that epoch count.
Three fictitious canary sentences are added to the training set to audit memorisation.

\subsection{Ruler-guided preference optimisation}
\label{sec:pref}
% 中文要點:
% 每前文 8 候選;分數 = 校準 log-odds,封頂在真實回覆的 95 百分位(只封上界;
% 更正 2026-10-17:原稿稱「典型帶」並引 Meister et al. 2023,但那是解碼時逐 token、雙側的帶,只能寫「受其啟發」);
% 約束 = 離題門檻(bge-m3 只做這件事)、抄襲、用語外漏;
% 行為格 = 長度 × 貼圖 × 連發 × 語氣詞,chosen 要在真實回覆的格內且 > +δ,rejected < −δ
% (更正 2026-10-17:分格 + 每格留最佳借自 MAP-Elites;「取真實回覆所在格」是我們加的規則,MAP-Elites 沒有目標分布);
% 沒有 chosen 就用真實回覆(DITTO,≤ 50%)。
% DPO β 0.1、參考固定、以驗收尺提早停止,不用 KL 懲罰。先報 best-of-8。

\paragraph{Candidates.}
For each training context $c$ the SFT model draws $n = 8$ candidates (temperature 0.8, top-p 0.95).
Each candidate $y$ receives a calibrated log-odds score $s(y)$ from the author LM, capped at the 95th percentile of the scores of T's real replies.
% NOTE (更正 2026-10-17, PROBLEMS_AND_SOLUTIONS.md S26): the draft called the cap a "typicality band" and cited
% Meister et al. 2023 as if it were the same object. Their band is per token, two-sided and on one model's information
% content; ours is a one-sided cap on whole-reply author log-odds. The code clamps the score: min(raw, cap) in persona/pairs.py.
The cap is one-sided: a reply that scores as more T-like than 95\% of T's own replies gains nothing from the excess, on the view that such a reply is a caricature rather than an imitation.
The cap is inspired by locally typical sampling \citep{meister2023locally}, but the quantity differs.
Locally typical sampling keeps, at each decoding step, the tokens whose information content is close to the conditional entropy of one model, which is a two-sided criterion applied per token.
Our cap is an upper bound only, and it applies to the author log-odds (T against S) of a whole reply.

\paragraph{Constraints.}
A candidate can be \emph{chosen} only if it is on topic (its bge-m3 cosine to the real reply exceeds the 90th percentile of cosines between unrelated real replies), is not a verbatim copy (no common substring of 12 or more characters with any training target), and contains no simplified-Chinese or Mainland-term leakage (the OpenCC converter \citepalias{opencc} in s2twp mode leaves it unchanged).
% NOTE: no bib entry exists for bge-m3 or Qwen3.5 in our docs; add them before submission.
Semantic embeddings are used \emph{only} for this on-topic gate.
On our data they separate on-topic from off-topic replies (AUC 0.94), but they do not see style (within-person vs cross-person difference 0.018), and in our check they did not separate an affirmative from a negative reply: the cosine between \zh{可以啊}{\textit{k\v{e}y\v{\i} a}} (``sure'') and \zh{不行耶}{\textit{b\`{u} x\'{\i}ng ye}} (``no way'') is 0.77.
We therefore do not use them to score style.

\paragraph{Behaviour cells.}
Candidates are binned by four interpretable descriptors: length bucket, sticker or not, number of messages in the burst, and sentence-final particle or not.
The \emph{target cell} of a context is the cell of T's real reply.
The chosen reply is the highest-scoring qualifying candidate with $s > +\delta$ in the target cell.
If the target cell holds no qualifying candidate, the neighbouring cells (same sticker flag, length bucket $\pm 1$) are searched, then any cell, and the level used is recorded and reported; with eight candidates an exact cell hit is rare.
The rejected reply is the lowest-scoring candidate in any cell with $s < -\delta$, where $\delta = \log 3$ in calibrated log-odds.
% NOTE (更正 2026-10-17, PROBLEMS_AND_SOLUTIONS.md S25): the draft said this "quality-diversity construction (MAP-Elites)
% makes the chosen set follow T's empirical distribution over behaviours rather than collapsing onto the single safest
% reply, which is the known failure mode of preference optimisation on dialogue". MAP-Elites has no target distribution,
% and "the known failure mode" had no source; rewritten below.
The cell structure borrows the archive of MAP-Elites \citep{mouret2015illuminating}, in which user-chosen descriptors partition the space and the best solution found so far is kept in each cell.
MAP-Elites itself is an evolutionary search that mutates the stored elites in order to fill as many cells as possible with high-performing solutions, and it has no target distribution over cells; taking the chosen reply from the cell of T's real reply is our addition.
The intent is that the chosen set follows T's empirical distribution over behaviours, to the extent that exact or neighbouring cell hits occur, instead of concentrating on a single safe reply, the loss of diversity that diversity-aware preference methods such as DivPO \citep{lanchantin2025divpo} also address.
If no candidate qualifies as chosen, the real reply is used, as in DITTO \citep{shaikh2025ditto}, limited to half of the pairs; if no candidate qualifies as rejected, the context yields no pair.

\paragraph{Objective.}
We use standard DPO \citep{rafailov2023dpo} with $\beta = 0.1$ against the frozen SFT model as reference, for one epoch, with LoRA on top of the merged SFT weights.
There is no KL penalty beyond the implicit one; instead, training stops when the \emph{evaluation} rulers (\S\ref{sec:eval}), computed on a development sample at 100-step intervals, stop improving.
% NOTE: handover.md §6 item 5: DPO early stopping is currently evaluated only at the end of an epoch; persona/train_dpo.py
% accepts eval_every=100 but does not use it yet. The sentence above describes the design, not the current code.
This follows the observation that KL penalties did not prevent over-optimisation in reward-model fine-tuning \citep{gao2023scaling}, and the Goodhart principle that an optimised measure stops measuring: the training ruler and the acceptance rulers are different objects.

\paragraph{Best-of-$n$ before DPO.}
Before DPO we report best-of-8 selection with the same score and constraints, in the style of RAFT \citep{dong2023raft}.
It needs no training and indicates how far the ruler alone can move the outputs.

\subsection{Evaluation protocol}
\label{sec:eval}
% 中文要點:
% 驗收尺 = 測試期重訓的獨立分辨器(報 AUC,附底線與靈敏度)、NLL、NCD、特徵帶、JSD、離題率、銜接、抄襲、外漏、多樣性、金絲雀。
% 四個對照組(真實、打亂的真實、S、底模)把指標的座標定出來。
% 本人盲判用序貫停止(更正 2026-10-17:是「每題後看 Clopper–Pearson 95% 區間是否整個在 0.65 一側」的反覆區間規則,
% 不是 Wald 1945 的 SPRT;每題都看,實際錯誤率會高於名目的 5%)。
% 三條驗收線預先登記。

\paragraph{Independent discriminator.}
The acceptance ruler is not the training ruler.
On the held-out period we train a fresh character-n-gram discriminator between T's real replies and a system's replies to the same contexts (5-fold out-of-fold), a classifier two-sample test in the sense of \citet{lopezpaz2017revisiting}, and report its AUC on single turns (D-AUC$_1$) and on bags of 10 (D-AUC$_{10}$).
% NOTE: 0.885 is the single-turn AUC of the training-period author LM on the filtered test turns (Exp. 4; PIPELINE_SPEC.md).
% The control run of 2026-10-13 with the test-period discriminator measured real T vs real S at 0.76 (k = 10: 1.00) and
% real vs real at 0.46 (handover.md §4 item 8). Decide which value defines the sensitivity line.
Two controls fix the scale: the \emph{noise floor} (real vs real, expected 0.5) and the \emph{sensitivity} (real T vs real S, 0.885 single-turn after filtering).
A system is read as ``closer to T than S is'' when its D-AUC is below the sensitivity value, and as indistinguishable when its interval covers the floor.
Because the summed discriminator score drifts with bag size, bag-level results are reported as AUC or calibrated, not with a fixed threshold.
Folds are assigned by reply string, so an identical reply cannot be in a training fold and a test fold at once; the shuffled-real control, which must sit at 0.5 by construction, verifies this (a per-context split let the discriminator memorise strings and gave 0.08).

\paragraph{Other automatic measures.}
% NOTE: rarefaction still needs a source (PROBLEMS_AND_SOLUTIONS.md S27).
The other automatic measures are: held-out negative log-likelihood of T's real replies under the model (a fit measure that needs no auxiliary model); the NCD score; the number of the 16 features inside T's tolerance band and the Jensen--Shannon divergence \citep{lin1991divergence} between generated and real feature distributions; on-topic rate; a coherence score $\log p(y \mid c) - \log p(y)$ under the base model, after the MMI objective of \citet{li2016diversity}; copy rate; OpenCC leakage per 100 characters; diversity (distinct-2 \citep{li2016diversity} and the Hill number of order 1 \citep{hill1973diversity} after rarefaction to equal sample sizes); and canary recall.

\paragraph{Controls that define the axes.}
Each measure is also computed for the real replies (ceiling), for \emph{shuffled} real replies (T's replies reassigned to other contexts: T's own style with the wrong content, which separates style measures from content measures), for S's real replies (negative control) and for the untrained base model.

\paragraph{Human judgement.}
T performs a blind two-alternative test (real vs system reply for the same context, ``which one did I write?'') and a preference test (best fine-tuned vs best prompting system, ``which one is more like me?'').
% NOTE (更正 2026-10-17, PROBLEMS_AND_SOLUTIONS.md S17): the draft called this rule a sequential test citing Wald 1945
% (SPRT). The text below describes what persona/human_eval.py does. Code vs text:
% - sprt_stop(n_correct, n, p0=0.65, max_n=100) returns stop/'below' when the upper Clopper-Pearson 95% bound is < p0,
%   stop/'above' when the lower bound is > p0, and otherwise stops only at n >= max_n. Despite its name it is not an SPRT.
% - score() evaluates the rule once on the answers it is given; checking after each item is left to whoever runs the
%   session (PIPELINE_SPEC.md §6.2 says the check is made item by item, and also calls the rule "SPRT").
% - score() applies the same p0 = 0.65 to kind-B (preference) items, while acceptance line (3) uses 0.60 for preference,
%   so the preference test has no stopping threshold of its own.
% - make_items() takes a priority_sids list and puts those items first; computing where the rulers disagree
%   (query-by-committee) is not in human_eval.py.
For the blind test, after each answered item we compute the exact (Clopper--Pearson) 95\% interval for T's accuracy, and we stop as soon as the interval lies entirely below or entirely above 0.65, or after 100 items.
This is a repeated confidence-interval rule, not the sequential probability ratio test of \citet{wald1945sequential}, which fixes two simple hypotheses (for example accuracy 0.5 against 0.65) and error rates $\alpha$ and $\beta$ in advance and stops when the cumulative log-likelihood ratio crosses one of two boundaries.
Because the interval is inspected after each item, the probability of stopping on the wrong side of 0.65 is higher than the nominal 5\%; \tbd{estimate the actual error rate of this rule, or replace it by a Wald SPRT}.
Items where the automatic rulers disagree are placed at the front of the queue (query-by-committee).

\paragraph{Acceptance lines (pre-registered; may only be tightened after the baselines are seen).}
\begin{enumerate}
\item D-AUC$_1 \le 0.75$ and D-AUC$_{10} \le 0.90$, and a paired-bootstrap improvement over the best prompting baseline whose interval excludes 0.
\item A band count not below that of the real held-out replies minus one (the real replies sit inside 12 of the 15--16 defined bands because of period drift, so an absolute threshold would be stricter than the person), diversity at least 80\% of real, copy rate at most 2\%, and leakage not above the real replies.
\item T's blind accuracy at most 0.65 and preference for the fine-tuned system at least 0.60.
\end{enumerate}
Passing (1) and (2) but not (3) is reported as a failure of the rulers, not a success of the model.

\paragraph{Statistics.}
We use three training seeds per system; a paired bootstrap over test contexts (2,000 resamples) for each system comparison \citep{bergkirkpatrick2012empirical}; Clopper--Pearson intervals for single-turn accuracies; and a turn-level bootstrap for bag-level accuracies, because overlapping bags are not independent.

\subsection{Privacy and consent}
\label{sec:privacy}
% 中文要點:
% 全程本機;去識別化;S 的訊息只當前文與對照作者,發表前要同意;權重不公開;只報統計與 ≤ 8 字的 n-gram。

All processing is local.
Speaker names are replaced by role labels; phone numbers, identity numbers, addresses and e-mail addresses are replaced by placeholders with hand-written Taiwan-specific patterns.
S's messages are used only as context and as the contrast author for ruler validation; S's consent is required before any publication that reports statistics derived from S's messages.
Adapter weights are treated as derived personal data and are not released.
The paper reports only aggregate statistics and n-grams of at most eight characters.

% experiments.tex -- Sections 4 (Experiments) and 5 (Limitations), converted from METHOD_NEURIPS.md §4--5 on 2026-10-17.
% One sentence per line. Each subsection's 中文要點 is the comment block directly under its heading.
% "% NOTE:" comments mark places where the wording differs from the Markdown draft, or where docs disagree.

\section{Experiments}
\label{sec:experiments}

\subsection{Tier 1: runnable now (one target person)}
\label{sec:tier1}
% 中文要點:
% 一人資料、約 700 訓練樣本;基準線 = zero-shot、many-shot(4 種 n)、檢索式、best-of-8;
% 主系統 = SFT(4 個資料量 × 2 底模)與 SFT+DPO;七項消融;六個預先登記假設;一張 16 GB 卡幾天內跑完。
% 預期:自動線贏 prompting、人判線離「分不出來」還遠。

\paragraph{Data.}
We use one LINE log (4,967 messages, 2025-06-20 to 2026-09-19).
% NOTE: the filtered counts below are Exp. 4 (loose rule; FORMAL_DISTANCE_EXPERIMENT.md §E). The stage-1 package with the
% emoticon whitelist gives train S 698 / T 707 and test S 118 / T 126 (handover.md §3; reports/ruler_validation.md).
% Pick one source before submission. reports/data_stats.json has the stage-1 sample counts for the TBD table.
The training period (before 2026-07-20) has 862 T turns (706 after symbol filtering) and 785 S turns (679); the held-out period has 149 T turns (126) and 134 S turns (115).
The median turn length is 16 characters.
Exact sample counts, context-length and target-length distributions, sticker share and per-month counts are reported in \tbd{Table: sample statistics from \texttt{samples.stats()}}.

\paragraph{Systems.}
Table~\ref{tab:systems} lists the systems and controls.
% NOTE: the draft's B1 purpose read "the strongest prompting alternative with the same data"; "strongest" is not yet
% measured, so the cell is rephrased.

\begin{table}[t]
\centering
\caption{Systems and controls in Tier~1.}
\label{tab:systems}
\begin{tabular}{@{}p{0.17\linewidth}p{0.43\linewidth}p{0.32\linewidth}@{}}
\toprule
id & system & purpose \\
\midrule
B0 & zero-shot, Qwen3.5-4B-Instruct, Chinese system prompt & prompting floor \\
B1-$n$ & many-shot ICL \citep{agarwal2024manyshot}, $n \in \{10, 50, 200, \text{all}\}$ example pairs from training & many-shot prompting with the same data; non-monotone in $n$ is expected \\
B2 & retrieval-augmented ICL, 20 nearest training contexts by bge-m3 & retrieval alternative \\
B3 & best-of-8 on SFT with the training ruler & ruler-only improvement, no preference training \\
SFT-$m$ & LoRA SFT on nested subsets $m \in \{100, 250, 500, \text{all}\}$, base and instruct & data-size curve; base vs instruct \\
SFT+DPO & \S\ref{sec:pref} on the better checkpoint, all data & main system \\
C-real / C-shuf / C-S / C-base & controls (\S\ref{sec:eval}) & define the axes \\
\bottomrule
\end{tabular}
\end{table}

\paragraph{Ablations.}
% NOTE: the draft's last ablation said "(to show over-optimisation)"; rephrased so the outcome is not presupposed.
The ablations use the best checkpoint, all data and 3 seeds: LoRA on attention projections only; a 3-hour segment gap; recency weighting (2026 samples $\times 2$); pairs without behaviour cells; no score cap; no real-reply fallback; and DPO with the training ruler used also for early stopping (to test for over-optimisation).

\paragraph{Pre-registered hypotheses.}
\begin{description}
\item[H1:] SFT-all beats each prompting system on D-AUC$_1$, NLL and the band count.
\item[H2:] The data-size curve is monotone from 100 to all on NLL but flattens on D-AUC (the discriminator saturates).
\item[H3:] SFT+DPO improves D-AUC$_{10}$ and band count over SFT-all without lowering the on-topic rate or diversity by more than 10\%.
\item[H4:] The base checkpoint yields higher diversity and lower D-AUC than instruct at equal data; the instruct checkpoint yields a higher on-topic rate.
\item[H5:] Removing the score cap or the behaviour cells raises the training score but worsens the band count and diversity (Goodhart).
\item[H6:] The blind accuracy of T on SFT+DPO is below that on the best prompting system.
\end{description}

\paragraph{Compute.}
% NOTE: handover.md §4 item 11 estimates 5--15 minutes for a full 3-epoch SFT run (10.1 GB VRAM, about 1 s/step);
% the 10--20 minutes below is the draft's estimate. Both are replaced by measured times.
Tier~1 uses one RTX 4070 Ti SUPER (16 GB).
We plan 24 SFT runs at an estimated 10--20 minutes each \tbd{measured SFT run time}, candidate generation of about 5,600 samples, and DPO runs of under an hour each \tbd{measured DPO run time}.
The whole tier is estimated to fit in a few GPU-days.

\paragraph{Expected outcome against prior work.}
% NOTE: the draft says Tier 1's contribution is "the validated-ruler methodology", while the main claim (main.tex
% comment) says the rulers are "an intermediate step, not the contribution"; the two should be reconciled.
IMPersona reports a 25\% pass rate against acquaintances with 500 private messages and 44\% with 13K \citep{shi2025impersona}.
Our 700 turns sit at the low end, so Tier~1 is expected to show a gap from ``indistinguishable'' on the human line while beating prompting on each automatic line.
The contribution of Tier~1 is the validated-ruler methodology and the data-size and base-vs-instruct curves for Traditional-Chinese short messages, which we did not find in the literature.

\subsection{Tier 2: planned (generalisation)}
\label{sec:tier2}
% 中文要點:
% 多人同意資料、多作者參照、變異數分解、公開語料只做尺的跨作者驗證、熟人盲測、貼圖編號。

\begin{itemize}
\item \textbf{Multiple consenting participants} (target 5--10, each with at least 500 turns, with partner identifiers in the context), to test whether the rulers' validity checks and the acceptance lines transfer across people, and whether group pre-adaptation followed by per-person LoRA, as in the personalised setting of \citet{ntt2024}, helps at this data size.
\item \textbf{Multi-author reference} for Cosine Delta and a generalisability-theory decomposition of feature variance into person $\times$ period $\times$ topic \citep{cronbach1972dependability}, which turns the validity check from pass/fail into an estimate of how many turns are needed.
\item \textbf{Public Taiwanese short-text corpora} (PTT comments, Dcard replies) for cross-author ruler validation only; they are not private messages and cannot replace Tier~1 data.
\item \textbf{Acquaintance blind tests} in place of self-judgement, following the IMPersona protocol \citep{shi2025impersona}, with a sequential stopping rule (\S\ref{sec:eval}).
\item \textbf{Sticker identity} once an export format with sticker ids is available.
\end{itemize}

\subsection{Baselines from the literature and why they are not directly comparable}
\label{sec:litbaselines}
% 中文要點:
% 最近的四個方法都沒有繁中短訊、單人幾百則、驗證過的尺;所以重現它們的訓練配方當基準線,不比論文數字。
% (英文正文改成「我們在這些論文裡沒找到」,不寫 None。)

IMPersona \citep{shi2025impersona} (Llama-3.1-8B, English iMessage, acquaintance Turing test), DITTO \citep{shaikh2025ditto} and DIAL \citep{zhu2025dial} (English, population-level or demonstration-level) and Flipping the Dialogue \citep{naous2026flipping} (user-simulator LMs) are the closest methods we found.
In these papers we did not find Traditional-Chinese short-message results, a single person at the few-hundred-turn scale, or rulers validated on the target pair.
We therefore reimplement their \emph{training recipes} as baselines (SFT with the IMPersona hyperparameters; DITTO pairs with the real reply as chosen; DIAL-style discriminator-ranked pairs without cells or cap) rather than comparing against their published numbers.

\section{Limitations}
\label{sec:limitations}
% 中文要點:
% 單人單對象、測試集小、過濾規則拿掉部分風格線索、對照作者只有一位、本人自判、實驗 4 未獨立重算。
% (更正 2026-10-17:單類別偵測器和負向選擇只能當類比,方法要引單類別分類文獻,出處待查。)

\begin{itemize}
\item There is one target person and one partner in Tier~1; partner accommodation effects (emoji and style converging to the interlocutor; \citealp{marko2022depends}) cannot be separated from T's style.
\item The held-out period is 126 turns; intervals are wide and all comparisons are paired.
\item The symbol filter removes 18\% of T's turns and 13.5\% of S's.
The turns flagged only by brackets or slashes are symmetric (about 31 per person), and the half-width punctuation cue from our analysis of 2026-10-03 turned out to be mostly a code artefact (on the filtered chat it is weak: bag-level AUC 0.68).
% NOTE (更正 2026-10-17, PROBLEMS_AND_SOLUTIONS.md S16): the draft cited Forrest et al. 1994 (negative selection) for the
% one-class detector. Negative selection generates random detectors and deletes those that match self samples; a
% one-class model of T shares only the use of self samples. Cited here as an analogy; the method needs a one-class reference.
\item The author LM is fitted with S as the only contrast author.
A one-class detector trained on T alone is planned but not validated.
Negative selection \citep{forrest1994self} is only an analogy for it: both need samples of ``self'' alone, but negative selection generates random detectors and deletes those that match self, which the planned detector does not do; \tbd{reference for one-class classification}.
\item Self-judgement by T is weaker evidence than judgement by acquaintances.
\item Exp.~4 (the filter decision) is a single run without independent reimplementation.
\end{itemize}


\bibliographystyle{plainnat}
\bibliography{refs}

\end{document}
